\documentclass[12pt,a4paper]{article}

\usepackage[margin=2.5cm]{geometry}
\usepackage{hyperref}
\usepackage{enumitem}
\usepackage{booktabs}

\hypersetup{
  colorlinks=true,
  linkcolor=blue,
  urlcolor=blue
}

\title{\textbf{Lilet\_r3ad} \\ \large AI-Powered Study Assistant --- Project Report}
\author{Project: lilet\_r3ad \\ Version: 0.1.0}
\date{November 2026}

\begin{document}

\maketitle
\tableofcontents
\newpage

%% ─────────────────────────────────────────────────────────────
\section{Overview}

Lilet\_r3ad is a browser-based AI study assistant that runs entirely on the client
side. Users upload course documents (PDF, DOCX, or plain text), and the application
automatically builds a structured outline, generates flashcard sets, creates
multiple-choice quizzes, and opens an AI tutor chat. The AI model is downloaded
once and cached in the browser via IndexedDB; every inference call runs locally
through WebGPU. Signed-in users can sync their book library to their configured
Supabase project.

Users sign in with email through Supabase Auth. Uploaded books are stored in a
Supabase \texttt{user\_books} table protected by row-level security; each account
can access only its own books. The active course, learning progress, and generated
content cache remain in browser \texttt{localStorage}, scoped to the signed-in user.

%% ─────────────────────────────────────────────────────────────
\section{Technology Stack}

\begin{table}[h]
\centering
\begin{tabular}{@{}ll@{}}
\toprule
\textbf{Layer} & \textbf{Technology} \\
\midrule
Language          & JavaScript (ES Modules, no transpilation) \\
Build tool        & Vite 8.3.6 \\
Runtime           & Browser (no Node.js at runtime) \\
AI inference      & @mlc-ai/web-llm 0.2.88 (WebGPU, Web Worker) \\
AI model          & \texttt{Qwen2.5-1.5B-Instruct-q4f16\_1-MLC} \\
Persistence       & Supabase (books and accounts); browser \texttt{localStorage} (active study state) \\
Styling           & Custom CSS (design tokens, no CSS framework) \\
Routing           & Hash-based router (custom, no library) \\
Testing           & Node.js built-in test runner (\texttt{node --test}) \\
Hosting           & Not specified in source code \\
\bottomrule
\end{tabular}
\caption{Full technology stack derived from \texttt{package.json} and source files.}
\end{table}

%% ─────────────────────────────────────────────────────────────
\section{AI Model}

The application uses \textbf{Qwen2.5-1.5B-Instruct-q4f16\_1-MLC}, a 1.5-billion
parameter instruction-tuned language model quantised to 4-bit floating point
(q4f16). The model identifier is defined as the constant \texttt{MODEL\_ID} in
\texttt{src/app/config.js} and consumed by \texttt{src/core/engine.js}.

Inference is executed in a dedicated Web Worker via
\texttt{CreateWebWorkerMLCEngine} from the \texttt{@mlc-ai/web-llm} package,
keeping the main thread unblocked during generation. A priority queue serialises
concurrent generation requests, giving user-initiated prompts priority over
background pre-generation tasks (flashcards, quizzes).

The model weights are fetched from the MLC CDN on first launch and stored
persistently in the browser's IndexedDB. The application calls
\texttt{navigator.storage.persist()} to prevent the browser from evicting cached
weights.

%% ─────────────────────────────────────────────────────────────
\section{User Journey}

\begin{enumerate}[leftmargin=*, label=\arabic*.]
  \item \textbf{Landing page} (\texttt{\#/}) --- The user sees the product hero,
        feature highlights, and can sign in or create an account.
  \item \textbf{Account access} (\texttt{\#/login}) --- Email/password sign-in
        and account creation through Supabase Auth.
  \item \textbf{Dashboard} (\texttt{\#/dashboard}) --- The user sees their
        document library. Clicking \emph{New document} opens a compact upload
        modal accepting PDF, DOCX, or TXT files.
  \item \textbf{Document processing} (\texttt{\#/app}) --- After upload, the app
        parses the document, splits it into sections using a structure-aware
        splitter (\texttt{src/document/}), and renders an interactive outline.
  \item \textbf{Mode selection} --- The user picks a study mode:
        \begin{itemize}
          \item \textbf{Teacher} (\texttt{\#/teacher}) --- AI-guided walkthrough
                with tutor chat per section.
          \item \textbf{Study Partner} (\texttt{\#/path}) --- Collaborative
                flashcard and quiz review.
        \end{itemize}
  \item \textbf{Final review} (\texttt{\#/final}) --- Summary of completed
        sections, quiz scores, and progress statistics.
\end{enumerate}

%% ─────────────────────────────────────────────────────────────
Book documents are synchronized through Supabase and protected by the row-level
security policy in \texttt{supabase/schema.sql}. Active course data, generated
flashcards, quiz results, and learning progress are stored in user-scoped browser
\texttt{localStorage}; configure \texttt{VITE\_SUPABASE\_URL} and
\texttt{VITE\_SUPABASE\_ANON\_KEY} for local and Vercel deployments.

%% ─────────────────────────────────────────────────────────────
\section{Limitations}

\begin{itemize}[leftmargin=*]
  \item \textbf{Local study state}: active course, generated content, and quiz
        progress remain in browser storage and do not yet synchronize between
        devices. The book library itself is cloud-synchronized.
  \item \textbf{Model load time}: the Qwen2.5 model (~1~GB) must be downloaded
        on first launch. On slow connections or low-end devices this takes several
        minutes.
  \item \textbf{WebGPU requirement}: the AI engine uses WebGPU, which is only
        available in desktop Chrome and Edge. Mobile browsers and Firefox are not
        supported for AI features.
  \item \textbf{Single-document study}: the application processes one document at
        a time; there is no multi-document or cross-document study mode.
  \item \textbf{Hosting}: no deployment configuration (Vercel, Netlify, Docker,
        etc.) was found in the source code.
\end{itemize}

\end{document}
